\documentclass[../main]{subfiles}
\begin{document}
\ifSubfilesClassLoaded{\setcounter{page}{4}}{}
\section{Il comunismo e l’individuo}
\label{sec:08_comunismo_individuo}

La questione dell’individuo costituisce uno dei temi cardine del pensiero comunista. In tale contesto, l’importanza attribuita all’individuo assume una risonanza inedita. Come affermano Marx ed Engels nel «Manifesto del Partito Comunista», «al posto della vetusta società borghese, con le sue classi e le contrapposizioni di classe, si afferma un’associazione in cui il libero sviluppo di ciascuno è la premessa fondamentale per il libero sviluppo di tutti»

L’intera vicenda passata deve essere interpretata non come una mera successione di eventi, bensì come un processo di formazione dell’essenza stessa dell’umanità. A questo punto, diviene evidente che si tratta anche di un percorso volto a plasmare un determinato tipo di personalità, in armonia con tale essenza. Quale debba essere, dunque, questa personalità? La risposta a questa domanda si concretizzerà di volta in volta in un contesto storico e pratico, ma in ogni caso si tratterà di una personalità capace di creare consapevolmente la storia all’interno di una comunità di individui simili, anch’essi dotati di spirito creativo. In una libera associazione di persone, l’intera esistenza di ciascuno sarà permeata dall’attività creativa, che rappresenta l’espressione più compiuta della personalità che la genera.

La personalità si configura come un fenomeno intrinsecamente sociale, che permea l'esistenza umana fin dalla sua genesi. L'individuo non emerge come tale alla nascita, bensì lo diviene attraverso un processo di integrazione all'interno delle dinamiche relazionali, assumendo un ruolo attivo nella società. Inizialmente, l'esperienza vitale del neonato è plasmata dall'esterno, dal contesto sociale in cui è immerso, con tutte le sfumature e le interazioni tra individui e tra l'uomo e l'ambiente naturale. Solo dopo aver interiorizzato, ovvero appropriato, le forme di espressione, o, più precisamente, i modelli culturali – che non sono in alcun modo innati nell'individuo – egli potrà sviluppare una vera e propria personalità. Questi modelli preesistono, pronti ad accogliere chiunque entri in scena. Per l'organismo biologico, essi rappresentano elementi esterni, ma ciò vale unicamente per l'organismo, non per la personalità in formazione. Il mondo dell'azione, inteso come l'ambito della cultura umana, costituisce il vero terreno di sviluppo della personalità. In una società comunista, ogni individuo deve appropriarsi di questo mondo, imparando a modificarlo e trasformarlo

Agendo, l’individuo comunista contribuirà in modo attivo alla trasformazione dei rapporti sociali, smettendo di esserne un mero prodotto passivo. L’intento è creare le condizioni affinché, nel corso del proprio sviluppo individuale, possa compiere nel più breve tempo possibile il percorso della maturazione sociale, proiettandosi in prima linea nella vita comunitaria. È imperativo ricercare e individuare incessantemente nuove forme di aggregazione che leghino l’individuo alla società, in modo tale da renderlo necessariamente un membro attivo e partecipe, capace di vivere appieno l’esistenza collettiva. Data la continua evoluzione della società, queste forme di «attività congiunta e differenziata» non possono essere stabilite in modo definitivo. La gestione dello sviluppo sociale, in una società comunista, consiste primariamente nell’apprendere a modificare tempestivamente i legami tra l’individuo e la società, in modo da indurlo ad agire con umanità in ogni specifico e nuovo contesto storico.

I comunisti sono spesso accusati di voler livellare le personalità, di aspirare a un’omologazione forzata degli individui. In questa critica, i difensori del capitalismo rivelano in realtà un aspetto peculiare della società capitalistica contemporanea. Nel sistema capitalistico, gli individui, nella loro generalità, tendono a uniformarsi a tal punto che la loro presunta individualità si riduce a una serie di dettagli superficiali e irrilevanti. Questa apparente individualità, tuttavia, è ben lontana dalla vera individualità, che si manifesta come espressione singolare della cultura umana universale interiorizzata da ciascun individuo. Il capitalismo, infatti, include gli individui in processi standardizzati di produzione e consumo, modellando il loro comportamento sia nell’ambito lavorativo che in quello privato. In sostanza, gli individui sono già, in larga misura, “uguali”. Il comunismo, al contrario, mira a creare uno spazio pubblico favorevole allo sviluppo dell’individualità.

I comunisti auspicano la creazione di una società in cui ogni individuo possa sviluppare appieno il proprio potenziale creativo. In una simile società, l'individualità di ciascuno si manifesterebbe come espressione attiva e realizzazione concreta del progresso sociale. In questo contesto, in che modo si può parlare di standard?

È chiaro che nessuno può aspirare a diventare un creatore universale senza aver prima acquisito un bagaglio culturale di base, essenziale per tale scopo. Di conseguenza, è necessario definire un insieme minimo di condizioni che la società deve garantire a ciascun suo membro, affinché questi possa sviluppare appieno le proprie capacità creative; deve quindi esistere uno standard minimo per tali condizioni. Quali sono gli elementi costitutivi di questo minimo?

Ciò implica un adeguato sviluppo del sistema sanitario, con particolare attenzione alla prevenzione e alla promozione di condizioni favorevoli a una vita sana per ciascun individuo, piuttosto che alla mera cura delle malattie; un livello e una qualità dell’istruzione garantiti a tutti; un sistema di ristorazione pubblica e, in generale, tutti i servizi pubblici in grado di migliorare la qualità della vita quotidiana; la protezione dell’infanzia e della terza età. In uno dei suoi scritti, Lenin affermò che il socialismo inizia con un bicchiere di latte per ogni bambino. È evidente che la semplice disponibilità di tale bicchiere di latte non implica un’evoluzione della società verso il comunismo, ma è innegabile che, senza garantire ai bambini un’adeguata alimentazione, la società non possa progredire in tale direzione. Analogamente, non può evolvere senza tutti gli elementi sopra menzionati (questa è ovviamente una lista incompleta; la sua definizione dovrebbe essere compito di esperti in diversi settori e, con il progresso della società, si amplierà).

Questo principio basilare di rispetto umano non può essere considerato un valore immutabile e definitivo. Esso si evolve in funzione delle circostanze in cui si verifica il progresso, in particolare in relazione al livello di maturazione raggiunto dalla cultura umana. Nell’attuale contesto, che dovrebbe fungere da stimolo per l’evoluzione della società verso un modello comunista, è imprescindibile includere l’istruzione superiore, la quale non debba limitarsi alla mera acquisizione di competenze informatiche di programmazione, ma che si estenda anche alle fondamenta scientifiche necessarie per comprendere l’insieme dei processi produttivi – un approccio politecnico accessibile a tutti. È altrettanto cruciale assicurare a tutti l’accesso ai moderni strumenti di comunicazione.

È doveroso sottolineare come persino le nazioni capitalistiche più evolute non possano esibire un tasso di alfabetizzazione universale né garantire l’accesso alle cure mediche all’intera popolazione.

Come già osservato, un livello minimo di umanità nei confronti dei singoli rappresenta un imperativo storico. Ma tale parametro minimo deve altresì contemplare un’espansione continua. Affinché una società comunista, nella quale il libero sviluppo di ogni individuo costituisca la condizione per lo sviluppo di tutti, possa esistere e prosperare, questo nucleo fondante deve costantemente incorporare i più recenti e significativi progressi compiuti dall’umanità. La società comunista sarà composta da individui intelligenti, sani e dotati, i quali, in ultima analisi, saranno gli unici in grado di organizzare la riproduzione sociale della propria esistenza su basi razionali.
\end{document}
